% !TeX root = closed-systems-from-comparison-completeness.tex
% ============================================================
\chapter{Quadratic Necessity, Four--Dimensional Diamonds, and Causality}
\label{ch:causality}
% ============================================================
\begin{chapterorientation}
\orientationitem{Input}{The stabilized quadratic carrier and the observer-side filtration inherited from the transport stack.}
\orientationitem{Role}{This chapter derives the causal-scaling consequences of a first visible quadratic obstruction.}
\orientationitem{Output}{Four-dimensional causal-diamond growth, Lorentzian cone structure in the realized branch, and the scalar channel laws later used by Hilbert reconstruction and macroscopic compatibility.}
\end{chapterorientation}
\section{Introduction}
Under the standing principle of closed-world admissibility
(\cref{sp:closed-world}), this chapter develops the causal-scaling
consequences of combining the stabilized quadratic carrier of
\cref{ch:interface} with the irreversible filtration of
\cref{ch:flow,ch:stitching}.  From those two inherited structures it
derives the macroscopic causal geometry used later in
\cref{ch:hilbert,ch:einstein}.
\Cref{ch:interface,ch:flow,ch:stitching} have already established two
structural facts of
the closed relational stack, and the present chapter begins from them
rather than re-proving them.
First, the first visible intrinsic obstruction carrier is quadratic.
At the intrinsic level, this carrier is the stabilized degree--\(2\)
layer, which we denote in this chapter by
\[
	\mathcal K:=\mathcal K_\infty \;\simeq\; F^2/F^3.
\]
At the realized level, the intrinsic--smooth interface developed in
\cref{ch:interface} shows that the second layer is forced intrinsically
and that, under faithful realization on this carrier, nonzero stabilized
square classes are equivalent to nonzero realized curvature.
Accordingly, before the present chapter begins, the stabilized quadratic
carrier already comes with the realized-level interface criterion used
here.  More precisely, the intrinsic stabilization is
\cref{thm:interface-seal}, and the curvature identification is
\cref{thm:interface-curvature}.
Second, the stack inherits the observer-side monotone filtration
arising from the irreversible observer descent of \cref{ch:flow}; in
\cref{ch:stitching} that filtration is shown to interleave
functorially with the refinement tower.  This inherited filtration
induces the causal preorder used below.
The purpose of the present chapter is to extract the macroscopic
consequences of these two structural facts.  More precisely, we prove
the following.
\begin{enumerate}[label=\textup{(\roman*)}]
	\item
	      Triangle closure leaves the quadratic layer unchanged.  Consequently the
	      first visible channel count is quadratic in the number of active
	      generators.
	\item
	      The observer-side filtration determines the correct
	      normalization at the first visible layer.  Together with the
	      quadratic channel count on the stabilized quadratic carrier,
	      that canonical normalization yields parabolic scaling and hence
	      four--dimensional causal-diamond growth.
	\item
	      Combined with the causal preorder, the first-variation pairing
	      extracted, under the inherited second-jet faithfulness condition,
	      from the realized degree--\(2\) channel attached to the
	      stabilized quadratic carrier is necessarily Lorentzian whenever
	      that pairing is nondegenerate.
\end{enumerate}
Thus, once the degree--\(2\) obstruction and the inherited
observer-side filtration are fixed, both the macroscopic diamond
dimension and the causal signature of the realized system are sharply
constrained.
\begin{remark}[Role of the chapter]
	The point of the chapter is not to re-establish from first principles
	that the first visible obstruction is quadratic.  That fact is already
	established in
	\cref{thm:interface-stabilized,thm:interface-seal}.  What must be
	proved here is subtler:
	every finite-generator triangle-closed model compatible with the stack
	inherits a quadratic visible carrier, and that quadraticity, together
	with the inherited observer-side filtration, has rigid macroscopic
	consequences.
\end{remark}
\begin{remark}[Why the number \(4\) is encoded one layer deeper]
	The exponent \(4\) does not first arise at the level of
	diamond-volume counting.  It is already determined by the interaction of two
	independent quadratic laws.
	The first visible graded piece is degree \(2\), so the visible channel
	count grows quadratically:
	\[
		Q(k)\asymp k^2.
	\]
	The filtration itself then determines linear accumulation of the first
	visible carrier:
	\[
		M(T)\asymp T.
	\]
	Hence
	\[
		k(T)\asymp \sqrt T.
	\]
	A causal diamond carries one filtration direction together with one
	quadratic visible carrier, so
	\[
		\mu(D_T)\asymp T\cdot Q(k(T))\asymp T^2.
	\]
	Passing from filtration time \(T\) to linear refinement scale
	\[
		L\asymp k(T)\asymp \sqrt T
	\]
	therefore yields
	\[
		\mu(D_L)\asymp L^4.
	\]
	Accordingly, the exponent \(4\) is the macroscopic shadow of the
	quadratic channel law together with the inherited observer-side
	irreversible filtration, whose canonical normalization makes
	refinement scale parabolic relative to filtration time.
\end{remark}
% ============================================================
\section{Triangle closure and the quadratic layer}
\label{sec:triangle-quadratic}
% ============================================================
Let
\[
	L_k=\langle e_1,\dots,e_k\rangle
\]
be the free local transport group on the \(k\) active first-order edge
germs.  Let
\[
	\varepsilon:\mathbb Z[L_k]\to \mathbb Z,
	\qquad
	J_k:=\ker(\varepsilon),
\]
and define the group-valued augmentation filtration
\[
	\mathfrak F^mL_k:=\{g\in L_k:g-1\in J_k^m\}.
\]
Thus the first visible generator module is
\[
	V_k:=\mathfrak F^1L_k/\mathfrak F^2L_k.
\]
\begin{lemma}[Degree--\(1\) carrier]
The classes
\[
	\bar e_1,\dots,\bar e_k,
	\qquad
	\bar e_r:=(e_r-1)\bmod J_k^2,
\]
form a basis of \(V_k\).  In particular,
\[
	V_k\cong \mathbb Z^k.
\]
\end{lemma}
\begin{proof}
This is the standard degree--\(1\) part of the Magnus expansion for the
free group on \(k\) generators: modulo \(J_k^2\), every word records only
the total signed exponent of each generator \cite{Magnus1935}.  Therefore the classes of the
\(k\) active edge germs form the indicated basis.
\end{proof}
\begin{lemma}[Quadratic commutator carrier]
\label{lem:quadratic-commutator-carrier}
The commutator pairing induces an isomorphism
\[
	\Lambda^2(V_k)
	\cong
	\mathfrak F^2L_k/\mathfrak F^3L_k.
\]
Under this isomorphism,
\[
	\bar e_i\wedge \bar e_j
	\longmapsto
	[e_i,e_j]\bmod \mathfrak F^3L_k.
\]
\end{lemma}
\begin{proof}
For \(a,b\in L_k\), the same augmentation calculation used in
\cref{lem:primitive-two-cell-magnus} gives
\[
	[a,b]-1\equiv (a-1)(b-1)-(b-1)(a-1)\pmod {J_k^3}.
\]
Thus commutator is alternating and bilinear on
\(\mathfrak F^1L_k/\mathfrak F^2L_k\), and it lands in
\(\mathfrak F^2L_k/\mathfrak F^3L_k\).  For a free group, the degree--\(2\)
part of the group-valued augmentation filtration is the degree--\(2\)
free Lie piece; its basis is given by the basic commutators
\([e_i,e_j]\), \(i<j\).  Hence the induced map from \(\Lambda^2V_k\) is an
isomorphism.
\end{proof}
\begin{lemma}[Primitive triangle boundary has degree--\(2\) leading term]
\label{lem:primitive-triangle-degree-two}
Let a primitive triangle witness have first-order edge classes
\(u,v,w\in V_k\) with boundary closure
\[
	u+v+w=0.
\]
If the triangle is strict, so that \(u\wedge v\neq0\), then its based
triangle holonomy has zero degree--\(1\) image and nonzero degree--\(2\)
image, equal up to orientation to \(u\wedge v\).
\end{lemma}
\begin{proof}
The degree--\(1\) term is the signed boundary sum \(u+v+w\), hence it
vanishes.  In the associated rational Lie algebra, the
Baker--Campbell--Hausdorff expansion of the ordered boundary product has
quadratic term
\[
	\frac12\bigl([w,v]+[w,u]+[v,u]\bigr).
\]
Using \(w=-u-v\) in the degree--\(1\) quotient, this expression reduces,
up to the sign determined by orientation, to a nonzero scalar multiple of
\(u\wedge v\).  Equivalently, in the integral Magnus filtration its
quadratic class is the corresponding primitive exterior class up to
orientation.  A strict triangle witness has nonzero area class
\(u\wedge v\neq0\), so its first nonzero group-filtration image is exactly
degree \(2\).
\end{proof}
\subsection{Interface seal inherited from the Stone-limit carrier}
The preceding computation explains why the first possible intrinsic
transport residue of a strict primitive triangle is quadratic.  The
Stone-limit statement used by the present chapter is the stronger
stabilized form proved in \cref{thm:interface-seal}: nontrivial triangle
obstruction cannot be supported purely in \(F^3\), and its first stable
nontrivial image is the quadratic carrier
\[
	\mathcal K\simeq F^2/F^3.
\]
Thus this chapter inherits, rather than assumes, the degree--\(2\)
carrier whose channel count and causal consequences are analyzed below.
% ============================================================
\section{Quadratic channel count}
\label{sec:quadratic-channel-count}
% ============================================================
\begin{theorem}[Quadratic channel count]
\label{thm:binomk2}
For \(k\) active independent first-order generators, the first visible
quadratic transport carrier has rank
\[
	Q(k):=\rank \Lambda^2(V_k)=\binom{k}{2}.
\]
In particular,
\[
	Q(k)\asymp k^2.
\]
\end{theorem}
\begin{proof}
By the degree--\(1\) carrier lemma, \(V_k\cong\mathbb Z^k\).  Therefore
\[
	\rank \Lambda^2(V_k)=\binom{k}{2}.
\]
By \cref{lem:quadratic-commutator-carrier}, this exterior square is exactly
the degree--\(2\) commutator carrier
\(\mathfrak F^2L_k/\mathfrak F^3L_k\), hence it is the first visible
quadratic transport carrier.
\end{proof}
\begin{remark}[Exterior-bilinear necessity]
The quadratic carrier is not merely quadratic in cardinality; it is
exterior-bilinear in structure:
\[
	\mathfrak F^2L_k/\mathfrak F^3L_k\cong \Lambda^2(V_k).
\]
The number \(\binom{k}{2}\) is therefore the rank of primitive oriented
\(2\)-cell channels determined by the active first-order generator module.
\end{remark}
% ============================================================
\section{Filtration normalization and parabolic scaling}
\label{sec:parabolic}
% ============================================================
Let \(T\) denote the observer-side filtration parameter inherited
from the irreversible observer descent of \cref{ch:flow} and
interleaved with the refinement tower in \cref{ch:stitching}. Let
\[
	k(T)
\]
denote the activated refinement depth at observer-side filtration
time \(T\), and define the activated quadratic mass by
\[
	M(T):=Q(k(T)).
\]
\begin{theorem}[Determined first-visible normalization]
	\label{thm:forced-normalization}
	There is a canonical normalization of the observer-side filtration
	time, inherited from the irreversible observer descent of
	\cref{ch:flow} and interleaved with the refinement tower in
	\cref{ch:stitching}, unique up to an overall positive
	multiplicative constant, for which
	\[
		M(T)\asymp T.
	\]
	Equivalently, this inherited filtration time is the cumulative count
	of independent first-visible degree--\(2\) transport classes.
\end{theorem}
\begin{proof}
	By \cref{thm:binomk2}, the degree--\(2\) visible mass at refinement
	depth \(k\) is
	\[
		Q(k)=\binom{k}{2},
	\]
	so for the activated depth \(k(T)\) one has
	\[
		M(T)=Q(k(T)).
	\]
	Define the normalized observer-side filtration parameter by
	cumulative visible mass:
	\[
		\tau(T):=M(T).
	\]
	This canonically normalizes the observer-side filtration inherited
	from the irreversible observer descent of \cref{ch:flow} and
	interleaved with the refinement tower in \cref{ch:stitching}, and
	any other admissible normalization of this inherited filtration
	differs by multiplication by a positive constant.
	With this choice,
	\[
		M(\tau)=\tau,
	\]
	hence, after renaming \(\tau\) as observer-side filtration time,
	\[
		M(T)\asymp T.
	\]
	Since \(Q(k)\) counts independent first-visible degree--\(2\)
	classes, the same normalization states equivalently that this
	observer-side filtration time is the cumulative count of these
	classes.
\end{proof}
\begin{lemma}[Parabolic scaling]
	\label{lem:parabolic}
	Under the canonical normalization of \cref{thm:forced-normalization},
	\[
		k(T)\asymp \sqrt T.
	\]
\end{lemma}
\begin{proof}
	By \cref{thm:binomk2},
	\[
		Q(k)\asymp k^2.
	\]
	By \cref{thm:forced-normalization},
	\[
		M(T)=Q(k(T))\asymp T.
	\]
	Therefore
	\[
		k(T)^2\asymp T,
	\]
	and hence
	\[
		k(T)\asymp \sqrt T.
	\]
\end{proof}
\begin{remark}[Intrinsic origin of the temporal parameter]
	The parameter \(T\) is not an externally imposed time coordinate.  It
	is the canonical normalization variable for the observer-side
	monotone filtration inherited from the irreversible observer descent
	of \cref{ch:flow} and interleaved with the refinement tower in
	\cref{ch:stitching}.

	This inherited filtration records the observer-side irreversible
	information loss carried across the refinement tower, and the causal
	preorder introduced below is induced by this filtration structure.
	Its normalization is fixed internally by the compatibility of two
	independent growth laws:
	\[
		M(T)\asymp T,
		\qquad
		Q(k)\asymp k^2.
	\]
	Compatibility determines
	\[
		k(T)\asymp \sqrt T.
	\]
	Thus \(T\) is determined, up to overall scale, as the filtration
	coordinate that linearizes irreversible coarse-graining relative to the
	quadratic refinement carrier.
\end{remark}
% ============================================================
\section{Diamond growth from closure and rectangular factorization}
\label{sec:diamonds}
% ============================================================
\subsection{Causal preorder}
Let \(\Phys\) denote the observational quotient.
\begin{definition}[Filtration preorder]
	For \(p,q\in\Phys\), define
	\[
		p\preceq q
	\]
	if at every observer-side filtration time inherited from the
	irreversible observer descent of \cref{ch:flow} and interleaved
	with the refinement tower in \cref{ch:stitching}, the image of
	\(p\) lies in the forward fiber of the image of \(q\).
\end{definition}
\begin{lemma}[The filtration preorder]
	The relation \(\preceq\) is a preorder.  If backward-fiber separation
	holds, then it descends to a partial order on the observational
	quotient.
\end{lemma}
\begin{proof}
	Reflexivity is immediate.
	For transitivity, suppose \(p\preceq q\) and \(q\preceq r\).  Then at
	every observer-side filtration time inherited from the irreversible
	observer descent of \cref{ch:flow} and interleaved with the
	refinement tower in \cref{ch:stitching}, the forward fiber of the
	image of \(p\) is contained in that of \(q\), and that of \(q\) is
	contained in that of \(r\).  Hence the forward fiber of \(p\) is
	contained in that of \(r\), so \(p\preceq r\).
	If both \(p\preceq q\) and \(q\preceq p\), then these observer-side
	forward fibers coincide at every such filtration time.  Under
	backward-fiber separation for that same inherited observer-side
	filtration, this determines observational equivalence.  Thus
	antisymmetry holds after quotienting by observational equivalence.
\end{proof}
For \(p\preceq q_T\), define the causal diamond
\[
	D(p,q_T):=\{z\in\Phys: p\preceq z\preceq q_T\}.
\]
% ------------------------------------------------------------
\subsection{Mixed-scale typing from closure}
\label{subsec:mixed-typing}
% ------------------------------------------------------------
We now remove the final free typing input from the diamond-growth
argument.
Recall from the closure theorem
\cref{thm:structural-closure} that
rectangular completeness determines the canonical product presentation
\[
	U \cong X_A\times X_B,
\]
the diagonal action of
\[
	G=\Aut(U,\mathcal C),
\]
the orbit quotient
\[
	\pi:X\to \Phys:=X/G,
\]
and, crucially, canonical descent: admissible state reports are exactly the
coherent ones, equivalently exactly the maps factoring through \(\pi\).
Recall also the structural classification of quotient extensions: any
non-endpoint dependence beyond quotient semantics is exhausted by
exactly two and only two mechanisms:
\begin{enumerate}[label=\textup{(\arabic*)}]
	\item representative selection via a section \(\Phys\to X\);
	\item morphism-level transport data not determined by endpoints.
\end{enumerate}
In the closed system treated here, representative selection is not part
of the admissible semantics.  Accordingly, at any finite mixed scale
\((k,T)\), the only admissible extra content is transport content.  At
that scale the transport content decomposes into exactly two visible
directions: the refinement-visible direction at depth \(k\), and the
observer-side-filtration-visible direction at time \(T\), where \(T\)
is measured in the observer-side filtration inherited from the
irreversible observer descent of \cref{ch:flow} and interleaved with
the refinement tower in \cref{ch:stitching}.  This is the precise
input needed for the mixed-scale typing theorem.
\begin{theorem}[Mixed-scale typing theorem]
	\label{thm:mixed-typing}
	Fix a refinement depth \(k\) and observer-side filtration time \(T\),
	inherited from the irreversible observer descent of
	\cref{ch:flow} and interleaved with the refinement tower in
	\cref{ch:stitching}.  Let \(B_{k,T}\) denote the Boolean algebra of
	admissible events on the corresponding mixed \((k,T)\)-quotient of
	the closed stack.
	Let
	\[
		B_k^{\mathrm{ref}}\subseteq B_{k,T}
	\]
	be the Boolean subalgebra generated by events whose truth value is
	determined entirely by the stage-\(k\) refinement-visible carrier, and
	let
	\[
		B_T^{\mathrm{fil}}\subseteq B_{k,T}
	\]
	be the Boolean subalgebra generated by events whose truth value is
	determined entirely by the observer-side filtration data up to time
	\(T\).
	Then
	\[
		B_{k,T}=\Bool\!\bigl(B_k^{\mathrm{ref}}\cup B_T^{\mathrm{fil}}\bigr).
	\]
	Equivalently, every admissible mixed-scale event is a Boolean
	combination of refinement events and observer-side filtration
	events.
\end{theorem}
\begin{proof}
	By canonical descent in the closed system, admissible state reports factor
	through the quotient
	\[
		\pi:X\to \Phys.
	\]
	Hence there is no admissible state-level content arising from arbitrary
	representative choice in \(X\).
	Now consider an admissible event at mixed scale \((k,T)\), where \(T\)
	is measured in the observer-side filtration inherited from the
	irreversible observer descent of \cref{ch:flow} and interleaved with
	the refinement tower in \cref{ch:stitching}.  By definition, such an
	event is compatible with quotient semantics and functorial under that
	mixed refinement/observer-filtration structure.  If its value were
	not determined by quotient-level endpoint data alone, then by the
	classification of quotient extensions it would have to arise from one of
	the two and only two enrichment loci: representative selection or
	morphism-level transport data.
	The first locus is excluded in the present closed-system semantics:
	introducing a section
	\[
		\Phys\to X
	\]
	would amount to adding representative structure, whereas admissibility
	here is exactly quotient admissibility.
	Therefore every admissible mixed-scale event can depend only on the
	second locus, namely transport data not determined by endpoints.
	At scale \((k,T)\), the visible transport data has exactly two typed
	sources.
	First, there is the refinement-visible source: dependence on the active
	stage-\(k\) refinement carrier.  Events of this type generate
	\(B_k^{\mathrm{ref}}\).
	Second, there is the observer-filtration-visible source: dependence on
	the observer-side filtration order or forward-fiber data up to time
	\(T\).  Events of this type generate \(B_T^{\mathrm{fil}}\).
	Since there is no third enrichment locus in the closed system, every
	admissible mixed-scale event is obtained by combining these two typed
	sources.  Because admissible events form a Boolean algebra, this means
	precisely that
	\[
		B_{k,T}=\Bool\!\bigl(B_k^{\mathrm{ref}}\cup B_T^{\mathrm{fil}}\bigr).
	\]
\end{proof}
\begin{remark}[What closure contributes]
	\Cref{thm:mixed-typing} is genuinely a consequence of closure.  Closure
	does not merely force quotient semantics for states; it also removes
	representative selection as an admissible source of mixed-scale content.
	What remains is exhausted by transport, and at scale \((k,T)\) the
	transport data is typed exactly by the refinement-visible carrier
	together with the observer-side filtration inherited from
	\cref{ch:flow} and interleaved with the refinement tower in
	\cref{ch:stitching}.  Thus the mixed-scale Boolean algebra has no
	third independent generator locus.
\end{remark}
% ------------------------------------------------------------
\subsection{Rectangular splitting and derived coarse multiplicativity}
\label{subsec:derived-coarse}
% ------------------------------------------------------------
By \cref{thm:mixed-typing}, the mixed-scale admissible-event algebra
\(B_{k,T}\) is generated by the refinement-event and filtration-event
subalgebras.  The lower Boolean layer of the stack identifies, for any
such generated pair, rectangular completeness, rectangle-algebra
factorization, and bijectivity of the canonical marginal ultrafilter
map.  Applying that theorem to
\[
	\bigl(B_k^{\mathrm{ref}},\,B_T^{\mathrm{fil}}\bigr)\subseteq B_{k,T}
\]
yields the mixed-scale rectangular splitting
\[
	B_{k,T}\cong B_k^{\mathrm{ref}}\otimes B_T^{\mathrm{fil}},
\]
equivalently the bijection
\[
	\Sigma_{k,T}:UF(B_{k,T})
	\longrightarrow
	UF(B_k^{\mathrm{ref}})\times UF(B_T^{\mathrm{fil}}).
\]
\begin{definition}[Activated refinement carrier]
	Fix a forward observer-side filtration ray \(T\mapsto q_T\),
	inherited from the irreversible observer descent of \cref{ch:flow}
	and interleaved with the refinement tower in
	\cref{ch:stitching}.  Let \(k(T)\) be the activated refinement
	depth.  Under the mixed-scale splitting, define the activated
	visible refinement carrier by
	\[
		\Chan_{k(T)}
		:=
		\pi^{\mathrm{ref}}\!\bigl(D(p,q_T)\bigr)
		\subseteq UF(B_{k(T)}^{\mathrm{ref}}),
	\]
	where \(\pi^{\mathrm{ref}}\) is the refinement projection.
\end{definition}
\begin{definition}[Product-limit measure]
	\label{def:product-measure}
	Fix compatible marginal measures
	\[
		\mu_k^{\mathrm{ref}}
		\quad\text{on}\quad UF(B_k^{\mathrm{ref}}),
		\qquad
		\mu_T^{\mathrm{fil}}
		\quad\text{on}\quad UF(B_T^{\mathrm{fil}}),
	\]
	and define a measure on \(UF(B_{k,T})\) by transporting the product
	measure through the mixed-scale ultrafilter splitting.  Assume that
	these measures are compatible under the bonding maps, so that they
	induce a limit measure \(\mu\) on the inverse limit.
\end{definition}
\begin{lemma}[Product comparability of diamonds]
	\label{lem:diamond-product}
	Assume \cref{def:product-measure}.  Let \(T\mapsto q_T\) be a
	forward observer-side filtration ray, inherited from the
	irreversible observer descent of \cref{ch:flow} and interleaved
	with the refinement tower in \cref{ch:stitching}, and let \(k(T)\)
	be the activated refinement depth at observer-side filtration time
	\(T\).
	Then there exist measurable sets
	\[
		S_{k(T)}\subseteq UF(B_{k(T)}^{\mathrm{ref}}),
		\qquad
		R_T,R_T'\subseteq UF(B_T^{\mathrm{fil}})
	\]
	such that
	\[
		S_{k(T)}\times R_T
		\;\subseteq\;
		\Sigma_{k(T),T}\bigl(D(p,q_T)\bigr)
		\;\subseteq\;
		S_{k(T)}\times R_T'
	\]
	and
	\[
		\mu_T^{\mathrm{fil}}(R_T)\asymp \mu_T^{\mathrm{fil}}([0,T]),
		\qquad
		\mu_T^{\mathrm{fil}}(R_T')\asymp \mu_T^{\mathrm{fil}}([0,T]),
	\]
	where the interval \([0,T]\) is measured in that observer-side
	filtration coordinate.
	Moreover one may take
	\[
		S_{k(T)}=\Chan_{k(T)}.
	\]
\end{lemma}
\begin{proof}
	Under the mixed-scale splitting for the observer-side filtration
	inherited from the irreversible observer descent of
	\cref{ch:flow} and interleaved with the refinement tower in
	\cref{ch:stitching},
	\[
		\Sigma_{k(T),T}:UF(B_{k(T),T})
		\longrightarrow
		UF(B_{k(T)}^{\mathrm{ref}})\times UF(B_T^{\mathrm{fil}})
	\]
	is a bijection.  Let
	\[
		\pi^{\mathrm{ref}},\ \pi^{\mathrm{fil}}
	\]
	denote the associated coordinate projections.
	Set
	\[
		S_{k(T)}:=\pi^{\mathrm{ref}}\!\bigl(D(p,q_T)\bigr)=\Chan_{k(T)}.
	\]
	By definition, every point of \(\Sigma_{k(T),T}(D(p,q_T))\) has
	refinement coordinate in \(S_{k(T)}\), so
	\[
		\Sigma_{k(T),T}(D(p,q_T))
		\subseteq
		S_{k(T)}\times \pi^{\mathrm{fil}}\!\bigl(D(p,q_T)\bigr).
	\]
	Define
	\[
		R_T':=\pi^{\mathrm{fil}}\!\bigl(D(p,q_T)\bigr).
	\]
	Because the preorder is induced by that observer-side monotone
	filtration, the observer-side filtration projection of a causal
	diamond is an interval in the observer-side filtration coordinate.
	More precisely, there exists an interval
	\(I_T\subseteq UF(B_T^{\mathrm{fil}})\) with
	\[
		I_T=\pi^{\mathrm{fil}}\!\bigl(D(p,q_T)\bigr)=R_T'.
	\]
	By the normalization of this inherited observer-side forward ray,
	this interval has observer-side filtration size comparable to the
	interval from \(0\) to \(T\):
	\[
		\mu_T^{\mathrm{fil}}(R_T')
		=
		\mu_T^{\mathrm{fil}}(I_T)
		\asymp
		\mu_T^{\mathrm{fil}}([0,T]).
	\]
	Now choose any measurable subinterval \(R_T\subseteq R_T'\) satisfying
	\[
		\mu_T^{\mathrm{fil}}(R_T)\asymp \mu_T^{\mathrm{fil}}(R_T')
		\asymp \mu_T^{\mathrm{fil}}([0,T]).
	\]
	Since \(S_{k(T)}\) is by definition the full refinement projection
	of the diamond, every refinement coordinate in \(S_{k(T)}\) occurs
	along that observer-side filtration ray, and shrinking only in the
	observer-side filtration coordinate yields
	\[
		S_{k(T)}\times R_T
		\subseteq
		\Sigma_{k(T),T}\bigl(D(p,q_T)\bigr)
		\subseteq
		S_{k(T)}\times R_T'.
	\]
	This proves the required product comparability.
\end{proof}
\begin{theorem}[Derived coarse multiplicativity]
	\label{thm:derived-coarse}
	Assume \cref{def:product-measure}.  Let \(T\mapsto q_T\) be a
	forward observer-side filtration ray, inherited from the
	irreversible observer descent of \cref{ch:flow} and interleaved
	with the refinement tower in \cref{ch:stitching}, and let \(k(T)\)
	be the activated refinement depth measured against that inherited
	observer-side filtration.  Then there exist constants \(c,C>0\)
	such that
	\[
		c\,\mu_T^{\mathrm{fil}}([0,T])\,
		\mu_{k(T)}^{\mathrm{ref}}(\Chan_{k(T)})
		\le
		\mu(D(p,q_T))
		\le
		C\,\mu_T^{\mathrm{fil}}([0,T])\,
		\mu_{k(T)}^{\mathrm{ref}}(\Chan_{k(T)}).
	\]
	Under the normalization conventions
	\[
		\mu_T^{\mathrm{fil}}([0,T])\asymp T,
		\qquad
		\mu_{k(T)}^{\mathrm{ref}}(\Chan_{k(T)})\asymp Q(k(T)),
	\]
	where \([0,T]\) is measured in that observer-side filtration, this
	yields
	\[
		\mu(D(p,q_T))\asymp T\,Q(k(T)).
	\]
\end{theorem}
\begin{proof}
	By the mixed-scale rectangular splitting for the observer-side
	filtration inherited from the irreversible observer descent of
	\cref{ch:flow} and interleaved with the refinement tower in
	\cref{ch:stitching},
	\[
		B_{k(T),T}\cong B_{k(T)}^{\mathrm{ref}}\otimes B_T^{\mathrm{fil}},
	\]
	equivalently the marginal ultrafilter map
	\[
		\Sigma_{k(T),T}:UF(B_{k(T),T})
		\longrightarrow
		UF(B_{k(T)}^{\mathrm{ref}})\times UF(B_T^{\mathrm{fil}})
	\]
	is bijective.
	By construction of the product-limit measure on this inherited
	observer-side filtration, \(\mu\) is the pushforward of
	\[
		\mu_{k(T)}^{\mathrm{ref}}\otimes \mu_T^{\mathrm{fil}}
	\]
	through \(\Sigma_{k(T),T}^{-1}\).
	By \cref{lem:diamond-product}, the image of the diamond is sandwiched
	between two product sets with identical refinement factor and
	comparable observer-side filtration factors:
	\[
		\Chan_{k(T)}\times R_T
		\subseteq
		\Sigma_{k(T),T}\bigl(D(p,q_T)\bigr)
		\subseteq
		\Chan_{k(T)}\times R_T',
	\]
	with
	\[
		\mu_T^{\mathrm{fil}}(R_T)\asymp \mu_T^{\mathrm{fil}}([0,T]),
		\qquad
		\mu_T^{\mathrm{fil}}(R_T')\asymp \mu_T^{\mathrm{fil}}([0,T]),
	\]
	where \([0,T]\) is taken in that observer-side filtration.
	Applying the product measure and transporting back through
	\(\Sigma_{k(T),T}^{-1}\) gives
	\[
		\mu_{k(T)}^{\mathrm{ref}}(\Chan_{k(T)})\,\mu_T^{\mathrm{fil}}(R_T)
		\le
		\mu(D(p,q_T))
		\le
		\mu_{k(T)}^{\mathrm{ref}}(\Chan_{k(T)})\,\mu_T^{\mathrm{fil}}(R_T').
	\]
	Using the comparability of \(R_T\) and \(R_T'\) with \([0,T]\) in
	that inherited observer-side filtration, we obtain constants
	\(c,C>0\) such that
	\[
		c\,\mu_T^{\mathrm{fil}}([0,T])\,
		\mu_{k(T)}^{\mathrm{ref}}(\Chan_{k(T)})
		\le
		\mu(D(p,q_T))
		\le
		C\,\mu_T^{\mathrm{fil}}([0,T])\,
		\mu_{k(T)}^{\mathrm{ref}}(\Chan_{k(T)}).
	\]
	Under the normalization conventions
	\[
		\mu_T^{\mathrm{fil}}([0,T])\asymp T,
		\qquad
		\mu_{k(T)}^{\mathrm{ref}}(\Chan_{k(T)})\asymp Q(k(T)),
	\]
	this becomes
	\[
		\mu(D(p,q_T))\asymp T\,Q(k(T)).
	\]
\end{proof}
\begin{definition}[Diamond dimension]
	For a forward ray \(T\mapsto q_T\) in the observer-side filtration
	inherited from the irreversible observer descent of
	\cref{ch:flow} and interleaved with the refinement tower in
	\cref{ch:stitching}, define
	\[
		\dim_\diamond
		:=
		\lim_{T\to\infty}
		\frac{\log \mu(D(p,q_T))}{\log k(T)},
	\]
	where \(k(T)\) is the activated refinement depth along that
	inherited observer-side filtration ray, whenever the limit exists.
\end{definition}
\begin{theorem}[Four-dimensional diamond growth]
	\label{thm:4D}
	Assume \cref{lem:parabolic,thm:derived-coarse}.  Let \(T\mapsto q_T\) be
	a forward ray in the observer-side filtration inherited from the
	irreversible observer descent of \cref{ch:flow} and interleaved with
	the refinement tower in \cref{ch:stitching}, and let \(k(T)\) be the
	corresponding activated refinement depth.  Then
	\[
		\mu(D(p,q_T))\asymp T^2.
	\]
	Equivalently, writing
	\[
		L:=k(T)\asymp \sqrt T,
	\]
	along that inherited observer-side filtration ray, one has
	\[
		\mu(D(p,q_T))\asymp L^4.
	\]
	In particular,
	\[
		\dim_\diamond=4.
	\]
\end{theorem}
\begin{proof}
	By \cref{thm:derived-coarse}, along a forward ray
	\(T\mapsto q_T\) in the observer-side filtration inherited from the
	irreversible observer descent of \cref{ch:flow} and interleaved with
	the refinement tower in \cref{ch:stitching}, with corresponding
	activated refinement depth \(k(T)\), one has
	\[
		\mu(D(p,q_T))\asymp T\,Q(k(T)).
	\]
	By \cref{thm:binomk2}, the quadratic channel count along that same
	inherited observer-side filtration ray satisfies
	\[
		Q(k(T))\asymp k(T)^2.
	\]
	By \cref{lem:parabolic}, the canonical normalization of that
	inherited observer-side filtration gives
	\[
		k(T)^2\asymp T.
	\]
	Therefore
	\[
		\mu(D(p,q_T))\asymp T\cdot T=T^2.
	\]
	Now set
	\[
		L:=k(T)\asymp \sqrt T.
	\]
	along that inherited observer-side filtration ray.
	Then
	\[
		T\asymp L^2,
	\]
	and consequently
	\[
		\mu(D(p,q_T))\asymp T^2\asymp (L^2)^2=L^4.
	\]
	Taking logarithms in the definition of \(\dim_\diamond\) along that
	same inherited observer-side filtration ray gives
	\[
		\dim_\diamond=4.
	\]
\end{proof}
\begin{remark}[Why \(4\) is determined]
	The exponent \(4\) records two inherited quadratic inputs: the
	quadratic channel law on the stabilized visible carrier, and the
	observer-side filtration whose canonical normalization makes
	refinement scale parabolic relative to filtration time.
	The first is algebraic:
	\[
		Q(k)\asymp k^2.
	\]
	The second quadraticity is kinematic only through the inherited
	observer-side filtration: its canonical normalization, interleaved
	with the refinement tower, forces
	\[
		T\asymp k(T)^2.
	\]
	Combining the quadratic channel law with that inherited
	observer-side normalization yields
	\[
		\mu(D_T)\asymp T\cdot Q(k(T))\asymp T\cdot T=T^2
	\]
	along the observer-side filtration, and therefore, after
	reparameterizing by the corresponding refinement scale \(L:=k(T)\),
	\[
		\mu(D_L)\asymp L^4.
	\]
	Thus the four-dimensional diamond law is the macroscopic shadow of the
	quadratic channel law together with the inherited observer-side
	irreversible filtration, whose canonical normalization makes
	refinement scale parabolic relative to filtration time.
\end{remark}
% ============================================================
\section{Lorentzian signature from the causal preorder}
\label{sec:lorentz}
% ============================================================
\Cref{sec:diamonds} identifies a \(4\)-dimensional first-variation
arena from the stabilized quadratic carrier together with the
inherited observer-side filtration and the macroscopic
diamond-scaling law they force.  The first-variation spaces themselves
belong to the reconstruction stack developed later in the manuscript;
what must be proved here is that the causal preorder already
determines the appropriate cone structure, and hence determines
Lorentzian signature on the first-variation pairing extracted, under
the inherited second-jet faithfulness condition, from the realized
degree--\(2\) channel attached to the stabilized quadratic carrier.
\begin{definition}[First-variation pairing]
	For \(\mu\)-almost every \(p\in S_\infty\), let
	\[
		\mathsf V_p
	\]
	denote the canonical \(4\)-dimensional first-variation space attached
	to \(p\), and let
	\[
		B_p:\mathsf V_p\times \mathsf V_p\to\mathbb R
	\]
	denote the symmetric bilinear first-variation pairing extracted,
	under the inherited second-jet faithfulness condition, from the
	realized degree--\(2\) channel attached to the stabilized quadratic
	carrier.  Write
	\[
		Q_p(v):=B_p(v,v).
	\]
\end{definition}
\begin{remark}[Logical status of \(\mathsf V_p\) and \(B_p\)]
	The existence of the first-variation spaces \(\mathsf V_p\) and of the
	pairings \(B_p\) is not an additional hypothesis of the present
	chapter.  It belongs to the later theorem-level reconstruction of the
	stack from the stabilized quadratic carrier together with, under the
	inherited second-jet faithfulness condition, the realized
	degree--\(2\) channel attached to that carrier.  The work of the
	present section is to derive, from the observer-side filtration
	preorder together with, under the inherited second-jet interface
	hypothesis, the causal organization of the negative directions of
	\(Q_p\) on the realized degree--\(2\) channel attached to the
	stabilized quadratic carrier.
\end{remark}
\subsection{Cone structure determined by the preorder}
\begin{definition}[Forward and backward first-variation cones]
	For \(\mu\)-almost every \(p\in S_\infty\), define
	\[
		\mathsf C_p^+\subseteq \mathsf V_p
	\]
	to be the set of first-variation directions represented by forward
	filtration transport at \(p\), and define
	\[
		\mathsf C_p^-:=-\mathsf C_p^+.
	\]
\end{definition}
\begin{theorem}[Causal cone compatibility]
	\label{thm:cone-compat}
	For \(\mu\)-almost every \(p\in S_\infty\), the preorder-induced future
	cone
	\[
		\mathsf C_p^+\subseteq \mathsf V_p
	\]
	satisfies the following properties:
	\begin{enumerate}[label=\textup{(C\arabic*)}]
		\item \label{C1}
		      \(\mathsf C_p^+\) is a proper convex cone containing no line;
		\item \label{C2}
		      for every \(v\in \mathsf C_p^+\),
		      \[
			      Q_p(v)\le 0;
		      \]
		\item \label{C3}
		      there exists \(v\in \mathsf C_p^+\) with
		      \[
			      Q_p(v)<0;
		      \]
		\item \label{C4}
		      \(\mathsf C_p^+\) contains two linearly independent vectors;
		\item \label{C5}
		      every vector with
		      \[
			      Q_p(v)<0
		      \]
		      lies in
		      \[
			      \mathsf C_p^+\cup \mathsf C_p^-,
		      \]
		      and
		      \[
			      \mathsf C_p^+\cap \mathsf C_p^-=\{0\}.
		      \]
	\end{enumerate}
\end{theorem}
\begin{proof}
	We prove the five statements in order.
	\smallskip
	\noindent
	\textit{Proof of \textup{\ref{C1}}.}
	By construction, \(\mathsf C_p^+\) consists of first-variation
	directions represented by forward filtration transport.  Since the
	filtration is monotone, composition of forward transport with forward
	transport is again forward.  Likewise, rescaling a forward infinitesimal
	variation by a nonnegative scalar preserves its forward orientation.
	Thus \(\mathsf C_p^+\) is a convex cone.
	It is proper because the filtration is irreversible: if a nonzero vector
	\(v\) and its negative \(-v\) both belonged to \(\mathsf C_p^+\), then
	the corresponding infinitesimal transport would be both forward and
	backward, determining a reversible local step and contradicting the arrow of
	coarsening established in
	\cref{thm:flow-inevitable-loss,cor:flow-irreversibility-forced}.  Hence
	\[
		\mathsf C_p^+\cap(-\mathsf C_p^+)=\{0\},
	\]
	so \(\mathsf C_p^+\) contains no line.
	\smallskip
	\noindent
	\textit{Proof of \textup{\ref{C2}}.}
	The intrinsic--smooth interface supplies, under the second-jet faithfulness condition, the realized-level criterion relating nonzero
	stabilized square classes to nonzero realized curvature, and the
	first-visible defect in the present causal setting is measured by the
	quadratic first-variation pairing.  Forward filtration transport is
	non-expansive at the first visible layer: along a forward variation,
	visible defect may persist or decrease, but it cannot become positive in
	the opposite orientation without violating monotonicity of the descent
	map.  Therefore every forward direction satisfies
	\[
		Q_p(v)\le 0
		\qquad
		(v\in \mathsf C_p^+).
	\]
	\smallskip
	\noindent
	\textit{Proof of \textup{\ref{C3}}.}
	If every forward direction satisfied \(Q_p(v)=0\), then, under the
	inherited second-jet faithfulness condition, the first-visible pairing
	on the realized degree--\(2\) channel attached to the stabilized
	quadratic carrier would vanish on the forward cone.  But the forward
	filtration is nontrivial, and the stabilized quadratic carrier is the
	first nonvanishing visible intrinsic obstruction layer.  Hence some
	forward direction must carry a nonzero first-visible class, and by
	\textup{\ref{C2}} this forces
	\[
		\exists\,v\in\mathsf C_p^+
		\quad\text{such that}\quad
		Q_p(v)<0.
	\]
	\smallskip
	\noindent
	\textit{Proof of \textup{\ref{C4}}.}
	By \cref{thm:4D}, causal diamonds have asymptotic growth law
	\[
		\mu(D_L)\asymp L^4.
	\]
	A cone generated by a single direction cannot support a genuinely
	four-dimensional first-variation arena: it would collapse the visible
	transport geometry to a one-parameter family and force a degenerate
	diamond law incompatible with \(L^4\)-growth.  Therefore the forward
	cone contains at least two linearly independent directions.
	\smallskip
	\noindent
	\textit{Proof of \textup{\ref{C5}}.}
	The two-locus classification theorem shows that all admissible
	non-endpoint transport data is exhausted by the transport locus.
	Within the transport locus, the monotone filtration supplies the only
	intrinsic orientation: forward or backward.  There is no third temporal
	orientation compatible with the closed-system comparison semantics.
	Let \(v\) satisfy
	\[
		Q_p(v)<0.
	\]
	Then \(v\) is a genuinely timelike first-visible direction.  By the
	exhaustion of admissible transport loci, \(v\) must be represented by
	transport that is either aligned with the filtration or opposite to it.
	Hence
	\[
		v\in \mathsf C_p^+\cup \mathsf C_p^-.
	\]
	Finally, if
	\[
		v\in \mathsf C_p^+\cap \mathsf C_p^-,
	\]
	then \(v\) is simultaneously forward and backward.  By irreversibility
	of the filtration, this is possible only for the zero direction.
	Therefore
	\[
		\mathsf C_p^+\cap \mathsf C_p^-=\{0\}.
	\]
	This proves \textup{\ref{C5}} and completes the theorem.
\end{proof}
\begin{remark}[Cone theorem as closure point]
	\Cref{thm:cone-compat} eliminates the last local hypothesis in the
	Lorentzian part of the chapter.  The cone structure is no longer
	assumed: it is determined by the observer-side filtration preorder
	together with, under the inherited second-jet faithfulness condition,
	the first-visible pairing on the realized degree--\(2\) channel
	attached to the stabilized quadratic carrier, the two-locus
	exhaustion theorem, and the four-dimensional diamond law already
	proved above.
\end{remark}
\subsection{Null directions and non-definiteness}
\begin{lemma}[Null directions exist]
	\label{lem:null-exists}
	Under the hypotheses of \cref{thm:cone-compat}, there exists
	\(w\neq 0\) such that
	\[
		Q_p(w)=0.
	\]
\end{lemma}
\begin{proof}
	By \textup{\ref{C3}}, choose
	\[
		v_0\in \mathsf C_p^+
	\]
	with
	\[
		Q_p(v_0)<0.
	\]
	By \textup{\ref{C4}}, choose
	\[
		v_1\in \mathsf C_p^+
	\]
	linearly independent from \(v_0\).  By \textup{\ref{C2}},
	\[
		Q_p(v_1)\le 0.
	\]
	If
	\[
		Q_p(v_1)=0,
	\]
	then \(w:=v_1\) is already a nonzero null vector and we are done.
	Hence we may assume
	\[
		Q_p(v_1)<0.
	\]
	Now consider the opposite direction
	\[
		-v_1\in \mathsf C_p^-.
	\]
	Since \(v_0\in \mathsf C_p^+\) and \(-v_1\in \mathsf C_p^-\), the two
	vectors lie in opposite timelike cones.  Consider the segment
	\[
		\gamma(s):=(1-s)v_0-sv_1,
		\qquad 0\le s\le 1.
	\]
	Its endpoints satisfy
	\[
		\gamma(0)=v_0\in\mathsf C_p^+,
		\qquad
		\gamma(1)=-v_1\in\mathsf C_p^-.
	\]
	We claim that the segment must leave the timelike set
	\[
		\{v:Q_p(v)<0\}.
	\]
	Indeed, if
	\[
		Q_p(\gamma(s))<0
		\qquad
		\text{for all }s\in[0,1],
	\]
	then by \textup{\ref{C5}} each \(\gamma(s)\) would belong to
	\[
		\mathsf C_p^+\cup\mathsf C_p^-.
	\]
	Because \([0,1]\) is connected and
	\[
		\mathsf C_p^+\cap \mathsf C_p^-=\{0\},
	\]
	a continuous path of timelike vectors joining \(v_0\in\mathsf C_p^+\)
	to \(-v_1\in\mathsf C_p^-\) would force a crossing from the forward cone
	to the backward cone inside the timelike region.  This is impossible
	without passing through the boundary separating the two sheets.
	Therefore there exists some \(s_0\in[0,1]\) such that
	\[
		Q_p(\gamma(s_0))=0.
	\]
	Set
	\[
		w:=\gamma(s_0).
	\]
	Since the endpoints of the segment are nonzero and lie in disjoint
	timelike cones, the boundary point separating the two sheets is
	nonzero.  Hence
	\[
		w\neq 0
		\qquad\text{and}\qquad
		Q_p(w)=0.
	\]
	This proves the existence of a nonzero null vector.
\end{proof}
\begin{proposition}[Non-definiteness determined by causality]
	\label{prop:non-definite}
	Under the hypotheses of \cref{thm:cone-compat}, the form \(B_p\) is not
	definite.  Equivalently, \(B_p\) is indefinite or degenerate.
\end{proposition}
\begin{proof}
	If \(B_p\) were positive definite or negative definite, then
	\[
		Q_p(v)\neq 0
		\qquad
		\text{for all }v\neq 0.
	\]
	This contradicts \cref{lem:null-exists}.  Therefore \(B_p\) cannot be
	definite.
\end{proof}
\subsection{Exclusion of split signature}
\begin{lemma}[Split signature is incompatible with two-sheeted cones]
	\label{lem:no-split}
	Let \(Q\) be a nondegenerate quadratic form on a real vector space \(V\)
	of dimension \(4\) with signature \((2,2)\).  Then the set
	\[
		\{v\in V:Q(v)<0\}
	\]
	is connected.  In particular, it cannot be written as a disjoint union
	\[
		\{Q<0\}=C^+\cup(-C^+),
		\qquad
		C^+\cap(-C^+)=\{0\},
	\]
	with \(C^+\) a proper convex cone.
\end{lemma}
\begin{proof}
	Choose coordinates on \(V\) in which
	\[
		Q(x)=x_1^2+x_2^2-x_3^2-x_4^2.
	\]
	Let
	\[
		v=(a,b,c,d),\qquad w=(a',b',c',d')
	\]
	be arbitrary vectors satisfying
	\[
		Q(v)<0,\qquad Q(w)<0.
	\]
	Then
	\[
		a^2+b^2<c^2+d^2,
		\qquad
		a'^2+b'^2<c'^2+d'^2.
	\]
	We first connect \(v\) to a vector whose first two coordinates vanish.
	Consider the path
	\[
		v(t):=(ta,tb,c,d),
		\qquad
		t\in[0,1].
	\]
	Along this path,
	\[
		Q(v(t))=t^2(a^2+b^2)-(c^2+d^2).
	\]
	Since \(0\le t^2\le 1\), we have
	\[
		t^2(a^2+b^2)\le a^2+b^2<c^2+d^2,
	\]
	hence
	\[
		Q(v(t))<0
		\qquad
		\text{for all }t\in[0,1].
	\]
	Thus \(v\) is joined within \(\{Q<0\}\) to the vector
	\[
		\tilde v:=(0,0,c,d).
	\]
	The same argument joins \(w\) within \(\{Q<0\}\) to
	\[
		\tilde w:=(0,0,c',d').
	\]
	It therefore suffices to show that the set
	\[
		E:=\{(0,0,u_3,u_4):u_3^2+u_4^2>0\}
	\]
	is connected.  But \(E\) identifies with
	\[
		\mathbb R^2\setminus\{0\},
	\]
	which is path connected.  Hence \(\tilde v\) and \(\tilde w\) are joined
	within \(E\), and therefore \(v\) and \(w\) are joined within
	\(\{Q<0\}\).
	Since \(v,w\in\{Q<0\}\) were arbitrary, the set \(\{Q<0\}\) is
	connected.
	Now suppose there were a decomposition
	\[
		\{Q<0\}=C^+\cup(-C^+),
		\qquad
		C^+\cap(-C^+)=\{0\},
	\]
	with \(C^+\) a proper convex cone.  Then \(C^+\) and \(-C^+\) would be
	disjoint nonempty relatively open subsets of \(\{Q<0\}\) whose union is
	all of \(\{Q<0\}\), contradicting connectedness.  Therefore no such
	decomposition exists.
\end{proof}
\begin{theorem}[Lorentzian signature is determined]
	\label{thm:lorentz-forced}
	Under the standing principle \cref{sp:closed-world}, if for
	\(\mu\)-almost every \(p\in S_\infty\) the first-variation pairing
	\(B_p\) is nondegenerate, then
	\[
		\sign(B_p)\in\{(1,3),(3,1)\}.
	\]
\end{theorem}
\begin{proof}
	A nondegenerate symmetric bilinear form on a \(4\)-dimensional real
	vector space has one of the signatures
	\[
		(4,0),\quad (3,1),\quad (2,2),\quad (1,3),\quad (0,4).
	\]
	By \cref{prop:non-definite}, the definite signatures \((4,0)\) and
	\((0,4)\) are excluded.
	By \cref{lem:no-split}, the split signature \((2,2)\) is excluded,
	because \cref{thm:cone-compat} gives a proper two-sheeted decomposition
	of the timelike set:
	\[
		\{Q_p<0\}=\mathsf C_p^+\cup \mathsf C_p^-,
		\qquad
		\mathsf C_p^+\cap \mathsf C_p^-=\{0\}.
	\]
	Hence the only remaining possibilities are
	\[
		(1,3)\quad\text{or}\quad(3,1).
	\]
\end{proof}
\begin{remark}[Status of the Lorentzian conclusion]
	The Lorentzian signature theorem is now conditional only on the
	nondegeneracy of the first-variation pairing itself.  The cone package
	is no longer assumed: it has been derived from the observer-side
	filtration preorder together with, under the inherited second-jet faithfulness condition, the first-visible pairing on the realized
	degree--\(2\) channel attached to the stabilized quadratic carrier.
	Thus, once nondegeneracy is granted, the causal organization needed
	for Lorentzian signature is already supplied by the observer-side
	filtration preorder together with, under the inherited second-jet faithfulness condition, the first-visible pairing on the realized
	degree--\(2\) channel attached to the stabilized quadratic carrier.
\end{remark}
% ============================================================
\section{Quadratic scalar laws and the Born rule}
\label{sec:born-law}
% ============================================================
The term ``Born rule'' is used here for the standard norm-square
probability law of Hilbert-space quantum mechanics \cite{VonNeumann1955}.
\Cref{thm:interface-stabilized,thm:interface-seal} identify the
stabilized quadratic carrier
\[
	\mathcal K \simeq F^2/F^3
\]
as the stabilized degree--\(2\) carrier through which the
first-visible defect data considered in this section are tracked in
the closed stack.  By \cref{thm:interface-seal}, a nontrivial
first-visible defect cannot hide purely in \(F^3\); accordingly the
first-visible alternatives considered here are tracked through their
classes in \(\mathcal K\).
In this section we derive the probability law for mutually exclusive
first-visible alternatives.  The derivation uses only the structural
constraints already established for admissible event algebras together
with the scalar-channel theory already proved for the realized
degree--\(2\) channel attached to the stabilized quadratic carrier
under the inherited second-jet faithfulness condition.
The argument has four steps.
\begin{enumerate}[label=\textup{(\roman*)}]
	\item
	      first-visible alternatives are tracked through their
	      first-visible classes in the stabilized degree--\(2\) carrier
	      \(\mathcal K\simeq F^2/F^3\);
	\item
	      admissible branch weights are scalarizations of the realized
	      degree--\(2\) channel attached to that stabilized quadratic carrier
	      under the inherited second-jet faithfulness condition;
	\item
	      the scalarization space of that realized channel is one-dimensional,
	      so every admissible branch weight is proportional to its canonical
	      scalar channel;
	\item
	      normalization removes the proportionality constant and yields the Born
	      law.
\end{enumerate}
% ------------------------------------------------------------
\subsection{First-visible alternatives}
% ------------------------------------------------------------
\begin{definition}[First-visible alternatives]
	A finite family of mutually exclusive realizable alternatives is called
	a family of \emph{first-visible alternatives} if their intrinsic
	first-visible defect classes are represented by elements
	\[
		v_1,\dots,v_n\in\mathcal K.
	\]
\end{definition}
These classes encode the first visible obstruction data associated with
the respective alternatives and, under the inherited second-jet
faithfulness condition, the probability-bearing scalar structure is
read through the realized degree--\(2\) channel attached to the
stabilized quadratic carrier.
% ------------------------------------------------------------
\subsection{Admissible branch weights}
% ------------------------------------------------------------
Let
\[
	\mathcal A=\{v_1,\dots,v_n\}\subset\mathcal K
\]
be a family of first-visible alternatives.
A probability assignment is a map
\[
	P:\mathcal A\to[0,1],
	\qquad
	\sum_{i=1}^n P(v_i)=1.
\]
We impose the following structural conditions.
\begin{enumerate}[label=\textup{(S\arabic*)}]
	\item \label{S1}
	      \textbf{Branch symmetry.}
	      Probabilities are invariant under permutations of alternatives with
	      identical intrinsic defect classes.
	\item \label{S2}
	      \textbf{Coarse-graining additivity.}
	      If two mutually exclusive alternatives are merged into a single coarse
	      alternative, then the probability of the coarse alternative is the sum
	      of the probabilities of the components.
	\item \label{S3}
	      \textbf{First-visible dependence.}
	      Probabilities depend only on the first-visible classes as
	      tracked, under the inherited second-jet faithfulness condition,
	      by the realized degree--\(2\) channel attached to the
	      stabilized quadratic carrier.
	\item \label{S4}
	      \textbf{Local order \(\le 2\), naturality, and linearity.}
	      The unnormalized branch-weight law is a carrier scalarization:
	      under the inherited second-jet faithfulness condition it depends
	      only on the realized degree--\(2\) channel attached to the
	      stabilized quadratic carrier, is natural under the realized
	      symmetry flow, and is linear on that realized channel.
\end{enumerate}
Condition \textup{(S2)} is determined by the Boolean algebra of admissible
events.  Condition \textup{(S3)} holds because no lower-order visible
carrier exists and because, under the inherited second-jet interface
hypothesis, the probability-bearing scalar structure is read only
through the realized degree--\(2\) channel attached to the
stabilized quadratic carrier.  Condition \textup{(S4)} is precisely
the scalar-channel structure already isolated for the realized
degree--\(2\) channel attached to the stabilized quadratic carrier
under the inherited second-jet faithfulness condition.
\begin{definition}[Admissible branch weight]
	\label{def:admissible-branch-weight}
	Under the inherited second-jet faithfulness condition, an
	\emph{admissible branch weight} is a map
	\[
		W:\mathcal K\to\mathbb R_{\ge 0}
	\]
	on the realized degree--\(2\) channel attached to the stabilized
	quadratic carrier, satisfying \textup{(S1)--(S4)}.
\end{definition}
% ------------------------------------------------------------
\subsection{Scalar reduction}
% ------------------------------------------------------------
\begin{lemma}[Scalar reduction]
	\label{lem:scalar-reduction}
	Under conditions \textup{(S1)--(S4)}, the probability law factors
	through an admissible branch weight \(W\) as
	\[
		P(v_i)=\frac{W(v_i)}{\sum_{j=1}^n W(v_j)}.
	\]
\end{lemma}
\begin{proof}
	By branch symmetry and first-visible dependence, probabilities
	depend only on the first-visible classes of the alternatives as
	tracked in \(\mathcal K\), while, under the inherited second-jet faithfulness condition, the probability-bearing scalar structure is
	read through the realized degree--\(2\) channel attached to the
	stabilized quadratic carrier.
	By coarse-graining additivity, the probability assigned to a coarse
	alternative is the sum of the probabilities of its mutually exclusive
	constituents.  Thus each alternative contributes a nonnegative scalar
	weight.
	By \textup{(S4)}, under the inherited second-jet interface
	hypothesis this scalar weight is an admissible scalarization on the
	realized degree--\(2\) channel attached to the stabilized quadratic
	carrier, hence an admissible branch weight in the sense of
	\cref{def:admissible-branch-weight}.
	Normalization then yields
	\[
		P(v_i)=\frac{W(v_i)}{\sum_j W(v_j)}.
	\]
\end{proof}
% ------------------------------------------------------------
\subsection{Uniqueness of the scalar channel}
% ------------------------------------------------------------
Let
\[
	Q:\mathcal K\to\mathbb R_{\ge0}
\]
denote a fixed nonzero canonical scalar channel on the realized
degree--\(2\) channel attached to the stabilized quadratic carrier
under the inherited second-jet faithfulness condition.
\begin{lemma}[Uniqueness of degree--\(2\) scalarization]
	\label{lem:K-scalar-unique}
	The space of admissible carrier scalarizations of the realized
	degree--\(2\) channel attached to the stabilized quadratic carrier
	under the inherited second-jet faithfulness condition is
	one-dimensional.  Equivalently, if
	\[
		W_1,W_2:\mathcal K\to\mathbb R
	\]
	are admissible carrier scalarizations, then there exists
	\[
		c\in\mathbb R
	\]
	such that
	\[
		W_1=c\,W_2.
	\]
\end{lemma}
\begin{proof}
	By the intrinsic--smooth interface, the scalar channel considered here
	is the realized degree--\(2\) channel attached to the stabilized
	quadratic carrier under the second-jet faithfulness condition.  In the
	scalar-channel theorem established later in the stack, the canonical
	scalar channel on that realized channel is proved to be well-defined
	up to an overall nonzero normalization constant because the scalar
	summand of the realized curvature module is one-dimensional.
	Equivalently, the space of admissible scalarizations of the realized
	degree--\(2\) channel attached to the stabilized quadratic carrier
	under the second-jet faithfulness condition is one-dimensional.
\end{proof}
\begin{lemma}[One-dimensionality of admissible branch weights]
	\label{lem:branch-weight-1d}
	The space of admissible branch weights is one-dimensional.
	Equivalently, every admissible branch weight \(W\) is of the form
	\[
		W=c\,Q
	\]
	for some constant \(c>0\).
\end{lemma}
\begin{proof}
	By \textup{(S4)}, under the inherited second-jet interface
	hypothesis an admissible branch weight is precisely a carrier
	scalarization of the realized degree--\(2\) channel attached to the
	stabilized quadratic carrier.  By \cref{lem:K-scalar-unique}, the
	space of scalarizations of that realized channel is one-dimensional.
	Hence if
	\[
		W:\mathcal K\to\mathbb R
	\]
	is any admissible branch weight and
	\[
		Q:\mathcal K\to\mathbb R
	\]
	is a fixed nonzero canonical scalar channel on that realized
	degree--\(2\) channel, then there exists
	\[
		c\in\mathbb R
	\]
	such that
	\[
		W=c\,Q.
	\]
	Since admissible branch weights are nonnegative and nontrivial on
	realizable alternatives, the proportionality constant must satisfy
	\(c>0\).
\end{proof}
% ------------------------------------------------------------
\subsection{Born law}
% ------------------------------------------------------------
\begin{theorem}[Born rule]
	\label{thm:born}
	Let
	\[
		v_1,\dots,v_n\in\mathcal K
	\]
	be first-visible alternatives.  Under conditions
	\textup{(S1)--(S4)}, the admissible probability law is
	\[
		P_i=\frac{Q(v_i)}{\sum_{j=1}^n Q(v_j)},
	\]
	where \(Q\) is the canonical scalar channel on the realized
	degree--\(2\) channel attached to the stabilized quadratic carrier
	under the inherited second-jet faithfulness condition.
\end{theorem}
\begin{proof}
	By \cref{lem:scalar-reduction}, probabilities arise from a normalized
	admissible branch weight \(W\):
	\[
		P_i=\frac{W(v_i)}{\sum_j W(v_j)}.
	\]
	By \cref{lem:branch-weight-1d},
	\[
		W=c\,Q
	\]
	for some \(c>0\).  Substituting gives
	\[
		P_i
		=
		\frac{c\,Q(v_i)}{\sum_j c\,Q(v_j)}
		=
		\frac{Q(v_i)}{\sum_j Q(v_j)}.
	\]
\end{proof}
\begin{remark}[Same source as the cosmological scalar]
	The scalar channel \(Q\) appearing in \cref{thm:born} is the same
	canonical degree--\(2\) scalar channel later isolated at the Einstein
	boundary on the realized degree--\(2\) channel attached to the
	stabilized quadratic carrier under the inherited second-jet interface
	hypothesis.  Accordingly, the Born probability law and the
	cosmological scalar term arise from the same scalar projection on that
	realized channel.
\end{remark}
% ============================================================
\section{Quadratic extension of the scalar channel}
\label{sec:quadratic-extension}
% ============================================================
Under the inherited second-jet faithfulness condition, the canonical
scalar channel on the realized degree--\(2\) channel attached to the
stabilized quadratic carrier extends to a quadratic form after real
linearization.  Once that quadratic extension is present, the bilinear
source is not an additional datum: it is forced by polarization.
\begin{definition}[Real linearized quadratic carrier]
	Under the inherited second-jet faithfulness condition, let
	\[
		K
	\]
	denote the realized degree--\(2\) channel attached to the stabilized
	quadratic carrier.  Define its real linearization by
	\[
		K_{\mathbb R}:=K\otimes_{\mathbb Z}\mathbb R.
	\]
\end{definition}
\begin{theorem}[Scalar polarization bridge]
	\label{thm:quadratic-extension}
	\label{thm:scalar-pairing}
	Under the inherited second-jet faithfulness condition, let
	\[
		Q_{\mathbb R}:K_{\mathbb R}\to\mathbb R_{\ge0}
	\]
	be the canonical quadratic extension of the degree--\(2\) scalar
	channel, with
	\[
		Q(v)=Q_{\mathbb R}(v\otimes 1)
		\qquad (v\in K).
	\]
	Then
	\[
		B_{\mathbb R}(x,y)
		:=
		\frac12\bigl(Q_{\mathbb R}(x+y)-Q_{\mathbb R}(x)-Q_{\mathbb R}(y)\bigr)
	\]
	is the unique symmetric bilinear form on
	\(K_{\mathbb R}\) whose diagonal is \(Q_{\mathbb R}\).  Its restriction
	\[
		B(v,w):=B_{\mathbb R}(v\otimes1,w\otimes1)
	\]
	is symmetric and biadditive on \(K\), and
	\[
		Q(v)=B(v,v)
		\qquad (v\in K).
	\]
\end{theorem}
\begin{proof}
	The map \(Q_{\mathbb R}\) is a quadratic form on the real vector space
	\(K_{\mathbb R}\).  Over a real vector space of characteristic not two,
	polarization assigns to any quadratic form a symmetric bilinear form by
	\[
		B_{\mathbb R}(x,y)
		=
		\frac12\bigl(Q_{\mathbb R}(x+y)-Q_{\mathbb R}(x)-Q_{\mathbb R}(y)\bigr).
	\]
	The quadratic identity for \(Q_{\mathbb R}\) gives additivity in each
	variable and homogeneity over \(\mathbb R\), so
	\(B_{\mathbb R}\) is bilinear.  Symmetry is immediate from the displayed
	formula.  Its diagonal is
	\[
		B_{\mathbb R}(x,x)
		=
		\frac12\bigl(Q_{\mathbb R}(2x)-2Q_{\mathbb R}(x)\bigr)
		=
		Q_{\mathbb R}(x),
	\]
	because \(Q_{\mathbb R}(2x)=4Q_{\mathbb R}(x)\).
	If \(C\) is any other symmetric bilinear form with
	\(C(x,x)=Q_{\mathbb R}(x)\), then
	\[
		C(x,y)=\frac12\bigl(C(x+y,x+y)-C(x,x)-C(y,y)\bigr),
	\]
	so \(C=B_{\mathbb R}\).  Thus the bilinear form is unique.
	Restricting \(B_{\mathbb R}\) to the embedded integral carrier
	\(K\hookrightarrow K_{\mathbb R}\) gives a symmetric biadditive pairing
	\(B:K\times K\to\mathbb R\).  Finally, for \(v\in K\),
	\[
		B(v,v)=B_{\mathbb R}(v\otimes1,v\otimes1)
		=Q_{\mathbb R}(v\otimes1)=Q(v).
	\]
\end{proof}
\begin{corollary}[Polarization identity]
	For the forced scalar-polarization bridge of \cref{thm:scalar-pairing}, the formula
	\[
		B_{\mathbb R}(x,y)
		=
		\frac12\bigl(Q_{\mathbb R}(x+y)-Q_{\mathbb R}(x)-Q_{\mathbb R}(y)\bigr)
	\]
	recovers the extended bilinear form.  Equivalently, because
	\(B_{\mathbb R}\) is symmetric,
	\[
		B_{\mathbb R}(x,y)
		=
		\frac14\bigl(Q_{\mathbb R}(x+y)-Q_{\mathbb R}(x-y)\bigr).
	\]
\end{corollary}
\begin{proof}
	The first identity is just the expansion formula for the quadratic form
	attached to a symmetric bilinear form:
	\[
		Q_{\mathbb R}(x+y)
		=
		Q_{\mathbb R}(x)+2B_{\mathbb R}(x,y)+Q_{\mathbb R}(y).
	\]
	Solving for \(B_{\mathbb R}(x,y)\) gives
	\[
		B_{\mathbb R}(x,y)
		=
		\frac12\bigl(Q_{\mathbb R}(x+y)-Q_{\mathbb R}(x)-Q_{\mathbb R}(y)\bigr).
	\]
	For the second identity, apply the first formula to \(x+y\) and \(x-y\):
	\[
		Q_{\mathbb R}(x+y)
		=
		Q_{\mathbb R}(x)+2B_{\mathbb R}(x,y)+Q_{\mathbb R}(y),
	\]
	\[
		Q_{\mathbb R}(x-y)
		=
		Q_{\mathbb R}(x)-2B_{\mathbb R}(x,y)+Q_{\mathbb R}(y).
	\]
	Subtracting yields
	\[
		Q_{\mathbb R}(x+y)-Q_{\mathbb R}(x-y)=4B_{\mathbb R}(x,y),
	\]
	hence
	\[
		B_{\mathbb R}(x,y)
		=
		\frac14\bigl(Q_{\mathbb R}(x+y)-Q_{\mathbb R}(x-y)\bigr).
	\]
\end{proof}
\begin{corollary}[Diagonal form]
	For the forced scalar-polarization bridge of \cref{thm:scalar-pairing}, there exists a basis of
	\(K_{\mathbb R}\) in which
	\[
		Q_{\mathbb R}(x)=\sum_{i=1}^r \lambda_i x_i^2,
		\qquad
		\lambda_i\ge 0.
	\]
\end{corollary}
\begin{proof}
	By \cref{thm:quadratic-extension}, \(Q_{\mathbb R}\) is represented by a
	symmetric bilinear form \(B_{\mathbb R}\) on the finite-dimensional
	real vector space \(K_{\mathbb R}\).  By the spectral theorem for
	symmetric bilinear forms over \(\mathbb R\), there exists a basis in
	which the representing matrix is diagonal:
	\[
		B_{\mathbb R}\sim \diag(\lambda_1,\dots,\lambda_r).
	\]
	For this standard finite-dimensional spectral fact, see
	\cite{Kato1995}.
	Since \(Q_{\mathbb R}(x)=B_{\mathbb R}(x,x)\) is nonnegative on the
	carrier by construction of the scalar channel, each diagonal entry
	satisfies
	\[
		\lambda_i\ge 0.
	\]
	Therefore in the corresponding coordinates,
	\[
		Q_{\mathbb R}(x)=\sum_{i=1}^r \lambda_i x_i^2.
	\]
\end{proof}
\begin{remark}[Why no amplitude language is introduced here]
	The present conclusion is purely real-linear.  It yields a diagonal
	quadratic form on \(K_{\mathbb R}\), but it does not by itself furnish a
	canonical complex structure.  Accordingly the chapter stops at the
	diagonal real form above.  Any later amplitude interpretation must be
	introduced separately and proved separately.
\end{remark}
% ============================================================
\section{Conclusion}
\label{sec:causality-conclusion}
% ============================================================
We summarize the logical output of the chapter.
First, triangle closure does not alter the quadratic graded layer.
The first visible transport carrier therefore remains the degree--\(2\)
carrier
\[
	\mathcal K \simeq F^2/F^3,
\]
and for a finite-generator triangle-closed model its visible channel
count is
\[
	Q(k)=\binom{k}{2}\asymp k^2.
\]
Second, the monotone filtration does not merely permit a normalization;
it determines one.  Under that canonical normalization, this inherited
filtration time is the cumulative count of independent first-visible
degree--\(2\) classes, so the cumulative first-visible mass is
canonically linear in observer-side filtration time:
\[
	M(T)\asymp T.
\]
Combining this with the quadratic channel law yields
\[
	k(T)\asymp \sqrt T.
\]
Third, closure and rectangular factorization force mixed-scale
splitting, and that splitting yields coarse multiplicativity of causal
diamonds:
\[
	\mu(D(p,q_T))\asymp T\,Q(k(T)).
\]
Substituting the quadratic channel law and the parabolic relation gives
\[
	\mu(D(p,q_T))\asymp T^2.
\]
Equivalently, in linear refinement scale
\[
	L\asymp \sqrt T,
\]
one obtains
\[
	\mu(D_L)\asymp L^4.
\]
Thus the macroscopic diamond dimension is determined to be \(4\).
Fourth, under the inherited second-jet faithfulness condition, the
observer-side filtration preorder canonically induces the forward
and backward causal cones in the first-variation space attached to
the realized degree--\(2\) channel of the stabilized quadratic
carrier.  Those cones are proper, convex, two-sided, and exhaustive
for timelike directions.
Accordingly the first-variation pairing extracted, under the
inherited second-jet faithfulness condition, from the realized
degree--\(2\) channel attached to the stabilized quadratic carrier
cannot be definite, cannot have split signature, and therefore, when
nondegenerate, must be Lorentzian:
\[
	\sign(B_p)\in\{(1,3),(3,1)\}.
\]
Finally, under the inherited second-jet faithfulness condition, the
realized degree--\(2\) channel attached to the stabilized quadratic
carrier controls the scalar sector.  Its admissible scalarizations
form a one-dimensional space, and this determines the probability
law for first-visible alternatives to be the normalized canonical
scalar channel on that realized channel:
\[
	P_i=\frac{Q(v_i)}{\sum_j Q(v_j)}.
\]
After real linearization of that realized degree--\(2\) channel under
the inherited second-jet faithfulness condition, the canonical scalar
channel extends to a nonnegative quadratic form on
\[
	K_{\mathbb R}=K\otimes_{\mathbb Z}\mathbb R,
\]
canonically diagonalizable over \(\mathbb R\).
The chapter therefore establishes a single rigid conclusion:
\begin{center}
	\emph{
		once the first visible obstruction is quadratic and the inherited
		observer-side filtration is fixed, the closed comparison stack
		determines quadratic channel count, parabolic refinement scaling,
		four-dimensional causal-diamond growth, and Lorentzian
		first-variation signature.
	}
\end{center}
No additional macroscopic dimension law and no alternative nondegenerate
causal signature are compatible with the closed relational architecture
developed up to this point once the stabilized quadratic carrier and
the inherited observer-side filtration are fixed.
\begin{remark}[Transition to the next chapter]
	The present chapter has extracted the large-scale causal consequences of
	the stabilized quadratic carrier together with the inherited
	observer-side filtration.  \Cref{ch:hilbert} follows the parallel
	linear branch: instead of deriving further causal scaling laws, it
	reconstructs Hilbert, phase, and scalar structure from the same
	stabilized quadratic layer.  This preserves the geometric and Hilbert
	realizations as parallel canonical outputs of one intrinsic quadratic
	obstruction package.
\end{remark}

\begin{chapterclosing}
\closingitem{Established}{The stabilized quadratic carrier and observer-side filtration determine the causal-diamond growth, parabolic scaling, and Lorentzian first-variation signature.}
\closingitem{Carried forward}{The parallel linear realization of the same carrier is developed next as Hilbert structure.}
\end{chapterclosing}
